\section{Frontend Architecture}
\label{sec:frontend-architecture}

\textbf{Status:} IMPLEMENTED.

\subsection{Route tree}
17 static routes (\texttt{/}, \texttt{about}, \texttt{blog}, \texttt{compiler}, \texttt{compare} + article, \texttt{contact}, \texttt{docs}, \texttt{features}, \texttt{languages}, \texttt{pricing}, \texttt{privacy}, \texttt{terms}, \texttt{use-cases} + 3), 2 dynamic segments (\texttt{languages/[slug]} $\times$12, \texttt{blog/[slug]} $\times$4), shared \texttt{layout.tsx} (navbar, footer, skip link), \texttt{not-found.tsx}, and metadata routes (\texttt{robots.ts}, \texttt{sitemap.ts}, \texttt{manifest.ts}). Route types are generated via \texttt{next typegen} (\texttt{LayoutProps<"/">} in \texttt{layout.tsx} reflects Next.js 16 conventions).

\subsection{Component catalogue}
\begin{table}[H]
\centering
\caption{Components (exhaustive)}
\begin{tabular}{|l|p{8.5cm}|}
\hline
\textbf{Component} & \textbf{Responsibility} \\ \hline
\texttt{RootLayout} (\texttt{layout.tsx}) & html lang, Inter/Plex Mono variables, full \texttt{metadata} (title template, canonical, OG/Twitter, robots), skip link, navbar/footer shell \\ \hline
\texttt{Navbar} & Logo (\texttt{next/image} favicon + wordmark + brand dot), route links with active state, mobile menu (\texttt{"use client"}) \\ \hline
\texttt{Footer} & 6-column link grid, copyright + brand lockup \\ \hline
\texttt{CodeEditor} & Window header (dots, title, status) + responsive iframe body; URL builder \\ \hline
\texttt{CompilerView} & Runtime switcher pills + remounted \texttt{CodeEditor} per language (\texttt{"use client"}) \\ \hline
\texttt{LanguageExplorer} & Search filter over the 12 pinned languages (\texttt{"use client"}) \\ \hline
\texttt{Breadcrumbs} + \texttt{JsonLd} & Visible trail plus \texttt{BreadcrumbList} schema; generic JSON-LD injector \\ \hline
\end{tabular}
\end{table}

\subsection{Props contract}
\texttt{CodeEditorProps} exposes \texttt{title}, \texttt{statusText}, \texttt{id}, \texttt{className}, \texttt{src}, \texttt{language}, \texttt{codeId}, \texttt{theme}, \texttt{fontSize}, and 12 embed flags. Defaults reproduce the legacy prototype URL (\texttt{theme=dark}, \texttt{hideTitle}, \texttt{codeChangeEvent}). Note an inconsistency documented in Section~\ref{sec:known-issues}: \texttt{buildOneCompilerSrc} defaults \texttt{availableLanguages} and \texttt{listenToEvents} to \texttt{true} while the component defaults them to \texttt{false}, so direct helper calls and default component renders produce different query strings.

\subsection{Styling approach}
Utility-first Tailwind classes over \texttt{@theme} tokens (\texttt{bg-canvas}, \texttt{bg-brand}, \texttt{text-ash}, \texttt{rounded-marketing}, \texttt{rounded-app-*}); no CSS modules. Icons are \texttt{lucide-react} only (\texttt{ArrowRight}/\texttt{ArrowDown}, \texttt{aria-hidden}, sized to text). \texttt{next/image} serves the favicon logo.

\subsection{State and data fetching}
Two \texttt{useState} islands only (\texttt{CompilerView} language, \texttt{LanguageExplorer} query). No context, no fetch, no client cache beyond framework defaults in \texttt{next.config.ts} (\texttt{cacheComponents}, partial prefetching).
